\documentclass[11pt]{article}
\usepackage[a4paper,margin=27mm]{geometry}
\usepackage{amsmath,amssymb,amsthm,xcolor,booktabs,array}
\newcommand{\PROVED}{\textcolor{green!40!black}{\textbf{PROVED}}}
\newcommand{\OPEN}{\textcolor{red!70!black}{\textbf{OPEN}}}
\newcommand{\AUDIT}{\textcolor{orange!70!black}{\textbf{AUDIT}}}
\newtheorem{theorem}{Theorem}
\newtheorem{corollary}{Corollary}
\title{Defect factorization and strict large-threshold induction (v113)}
\date{13 August 2026}

\begin{document}
\maketitle

\section{The exact defect identity}

Let $C:\mathscr H\to\mathscr K$ be a contraction and put
\begin{equation}
 D=I-C^*C,qquad D_*=I-CC^*.                                  \tag{1}
\end{equation}
In the normalized threshold notation let $h$ be the harmonic reference
lift and let $y$ be the opposite newborn feature.  The normalized cross is
\begin{equation}
 q=h-C^*y.                                                     \tag{2}
\end{equation}

\begin{theorem}[Julia defect reduction --- \PROVED]
The return load and its normalized inverse are
\begin{align}
 u:=Dh-q&=C^*(y-Ch),                                          \tag{3}\\
 D^{\dagger/2}u&=C^*D_*^{\dagger/2}(y-Ch),                    \tag{4}
\end{align}
on the natural closed ranges.  In particular, if the residual innovation
has a source factorization
\begin{equation}
                         y-Ch=D_*^{1/2}z,                      \tag{5}
\end{equation}
then
\begin{equation}
                         \|D^{\dagger/2}u\|\le\|z\|.          \tag{6}
\end{equation}
\end{theorem}

\begin{proof}
Substituting (2) gives
\[
 Dh-q=(I-C^*C)h-h+C^*y=C^*(y-Ch),
\]
which is (3).  Functional calculus applied to
$C^*D_*=DC^*$ gives the defect intertwining
$D^{1/2}C^*=C^*D_*^{1/2}$ and its Moore--Penrose version (4).  Under (5),
(4) becomes $D^{\dagger/2}u=C^*z$ after projection to the supported
range, proving (6).
\end{proof}

This identity explains the exact role of a lossless Julia network: it does
not need to invert the old signed operator.  It has only to produce the
innovation $z$ in (5).

\section{Arithmetic innovation}

The finite prime Julia colligation of G39 is lossless on the arithmetic
label space.  The delay synthesis of v111 is an isometry from that label
space to the two physical newborn delay families.  Their composition
therefore gives (5) for the arithmetic residual, with
\begin{equation}
 \|z_N^{\rm ar}(e)\|^2
 =s_N\|e\|^2,
 \qquad
 s_N:=\sum_{k\ge1}\frac{V_{N,k}}{(\log N)^{2k}}
 \longrightarrow s_*:=
 \frac{\sqrt\pi}{2}e^{1/4}\operatorname{erf}\!\left(\frac12\right).
                                                                    \tag{7}
\end{equation}

\begin{theorem}[Candidate physical arithmetic defect factorization --- \AUDIT]
If the G39 divisor-incidence output, after transport by the threshold delay
isometry of v111, is equal to the arithmetic part of the Kato-normalized
physical residual, then (5) holds with $z=z_N^{\rm ar}$ and hence
\begin{equation}
 \bigl\|D^{\dagger/2}u_N^{\rm ar}(e)\bigr\|^2
 \le s_N\|e\|^2,
 \qquad s_N\longrightarrow0.5922965364\ldots<1.               \tag{8}
\end{equation}
\end{theorem}

\begin{proof}
The prime Julia tensor product is unitary, so its output equation is the
standard defect equation (5) on the label innovation space.  Theorem 1 of
v111 identifies every label output with the corresponding physical delay
and preserves all inner products.  Equation (0) of v110 identifies its
depth-$k$ input with $(\log N)^{-k}\mathcal C_N^k\delta_1$, the G39
logarithmic Krylov word.  Under the stated equality with the physical
residual, its source norm is (7), and Theorem 1 gives (8).
\end{proof}

The hypothesis in this theorem is not supplied merely by equality of the
label and delay Grams.  It is the operator intertwining
\begin{equation}
 y_N^{\rm ar}-C_Nh_N^{\rm ar}
 =D_{*,N}^{1/2}\,
   \mathfrak U_N\mathfrak W_N,                                \tag{8a}
\end{equation}
where $\mathfrak U_N$ is the explicit delay isometry of v111.  G39 fixes
the divergence of the two sides, while (8a) asserts equality of the full
innovations.  The transfer-coefficient audit in v114 proves that the
\emph{arithmetic-only} form of (8a) is false: its zeroth-order part would
unitarily identify $J_{a,-}P_T$ with $J_{a,+}P_T$, although their primitive
Grams differ by $P_T(S_a+S_{-a})P_T\ne0$.  Thus Theorem 2 is retained only
as a conditional implication and cannot be used in the proof of row (d).
A possible replacement must mix the arithmetic ports with Gamma and Tate;
defining that replacement from the desired Schur return would be circular.

\section{Propagation with the Gamma--Tate remainder}

Suppose the old normalized defect has a fixed gap
\begin{equation}
                              D_N\ge\eta I.                    \tag{9}
\end{equation}
Then $\|C_N\|^2\le1-\eta$ and also $D_{*,N}\ge\eta I$ on the
opposite supported space.  By v112, the normalized Gamma--Tate residual
has norm $\varepsilon_N=O(N^{-1/4})$; it may therefore be written as
$D_{*,N}^{1/2}z_N^{\Gamma T}$ with
\begin{equation}
                    \|z_N^{\Gamma T}\|
 \le\eta^{-1/2}\varepsilon_N.                                \tag{10}
\end{equation}

\begin{corollary}[Strict large-threshold induction assuming (8a) --- \PROVED]
Let $0<\eta<(1-s_*)/4$.  There exists $N_0(\eta)$ such that, whenever
$N\ge N_0(\eta)$ and the old defect satisfies (9), the complete next-cell
capacity is at most $1-\eta$ and the next normalized defect again satisfies
(9).
\end{corollary}

\begin{proof}
Combine the arithmetic and Gamma--Tate source innovations.  Equations
(6)--(8) and (10) give
\begin{equation}
 \operatorname{Cap}_{N+1}^{1/2}
 \le\sqrt{s_N}+\eta^{-1/2}\varepsilon_N.                      \tag{11}
\end{equation}
Since $s_N\to s_*<1$ and $\varepsilon_N\to0$, the square of the right
side is at most $1-\eta$ for all sufficiently large $N$.  The normalized
Schur complement is $I$ minus this return capacity, so its defect is at
least $\eta I$.  This propagates (9).
\end{proof}

The corollary changes the global task substantially: after the exact
arithmetic transport, no new estimate is required at arbitrarily large
thresholds.  It remains to produce one certified base threshold
$N_b\ge N_0(\eta)$ with the normalized defect gap (9), including every
intermediate cell from the already certified $N=4$ endpoint.  This is a
finite, falsifiable interval-arithmetic task.  The existing
$T=\tfrac12\log5$ audit does not yet supply it because its eliminated safe
block was not coupled to the full complement.

\begin{center}
\begin{tabular}{>{\raggedright\arraybackslash}p{.55\linewidth}
                >{\centering\arraybackslash}p{.15\linewidth}
                >{\raggedright\arraybackslash}p{.20\linewidth}}
\toprule
Statement & Status & Location \\
\midrule
Exact Julia defect reduction & \PROVED & (1)--(6) \\
Strict norm of the candidate arithmetic innovation & \PROVED & v109--v111 \\
Arithmetic-only innovation intertwining (8a) & false & v114 \\
Cross-place Gamma--Tate replacement & \OPEN & independent conservation needed \\
Gamma--Tate perturbation under a fixed gap & \PROVED & v112, (9)--(10) \\
Large-threshold propagation assuming (8a) & \PROVED & Corollary 3 \\
Certified finite bridge from $N=4$ to one $N_b\ge N_0$ & \OPEN & interval computation required \\
Condition $G$, row (d) & \OPEN & not obtained by this conditional induction \\
\bottomrule
\end{tabular}
\end{center}

\end{document}
